% =====================================================================
%  Vorlage: Betriebsanweisung für ein Arbeitsmittel (Maschine, Gerät)
%  Kopieren nach betriebsanweisung/ba_<geraet>.tex und ausfüllen.
%  Nummer nach README_betriebsanweisung.md: BA-<Kürzel>-<Nr.>
%  Symbole: fablab-document/ba-symbole/QUELLEN.md
% =====================================================================
\BAsetup{
  typ=maschine,
  titel={Gerät Hersteller Modell (Kurzbeschreibung)},
  kurztitel={Gerät},                  % für das Inhaltsverzeichnis der Einweisung
  taetigkeit={Bedienen, Reinigen},
  nummer=BA-XX-01,
  stand=TT.MM.JJJJ,                   % Datum der Freigabe
  verantwortlich={Name der verantwortlichen Person},
  bild=bilder/geraet,                 % Foto, mehrere mit Komma
  entwurf,                            % entfernen, sobald freigegeben
}

\BAkopf

\BAblock{ANWENDUNGSBEREICH}{M002}{%
  \begin{itemize}
    \item Wofür das Gerät verwendet wird.
    \item Benutzung nur nach Einweisung; Hinweise des Herstellers beachten.
  \end{itemize}}

\BAblock{GEFAHREN FÜR MENSCH UND UMWELT}{W001}{%
  \begin{itemize}
    \item \textbf{Gefahr}: Beschreibung.
  \end{itemize}}

\BAblock{SCHUTZMASSNAHMEN UND VERHALTENSREGELN}{M004}{%
  \begin{itemize}
    \item Maßnahme.
  \end{itemize}}

\BAblock{VERHALTEN BEI STÖRUNGEN}{W001,F001}{%
  \begin{itemize}
    \item Gerät ausschalten, Betreuer informieren.
    \item Entstehungsbrand: Gerät stromlos machen, löschen. \textbf{Feuerwehr: 112}
  \end{itemize}}

\BAblock{VERHALTEN BEI UNFÄLLEN – ERSTE HILFE}{E003}{%
  \begin{itemize}
    \item Ruhe bewahren, Gerät abschalten, Betreuer/Ersthelfer verständigen.
    \item Jeden Unfall den Betreuern melden und im Verbandbuch eintragen.
  \end{itemize}\vspace{1.5mm}
  \textbf{Notruf: 112}\qquad FabLab: 09131\,/\,85\,28013}

\BAletzterblock{INSTANDHALTUNG, ENTSORGUNG}{M021}{%
  \begin{itemize}
    \item Wartung nur durch eingewiesene Betreuer im ausgeschalteten Zustand.
    \item Elektrische Prüfung nach DGUV Vorschrift 3 gemäß festgelegter Prüffrist.
  \end{itemize}}
